\documentclass[11pt]{article}
\usepackage[a4paper,margin=1in]{geometry}
\usepackage{hyperref}
\usepackage{longtable}
\usepackage{array}
\usepackage{xcolor}

\hypersetup{
  colorlinks=true,
  linkcolor=blue,
  urlcolor=blue,
  pdftitle={Agora: one sentence, one world},
  pdfauthor={Agora 2.0}
}

\title{\textbf{Agora: one sentence, one world}\\
\large Feature and Scope Technical Report}
\author{Agora 2.0 Runtime and Creator Pipeline}
\date{July 2, 2026}

\begin{document}
\maketitle

\begin{abstract}
Agora 2.0 converts a compact natural-language world brief into a package-backed, multiplayer pixel simulation. The current implementation is a durable workflow rather than a single blocking request: it creates a draft, queues generation, records revision artifacts, runs an art and QA worker, publishes a package, exposes that package through a strict catalog, and serves live sessions through REST and WebSocket APIs. The product promise is simple: one sentence can become one coherent world, but the engineering contract behind that promise is explicit, staged, validated, and database-backed.
\end{abstract}

\section{Product Thesis}

The phrase ``Agora: one sentence, one world'' describes the desired user experience: a creator provides a short scenario seed, and Agora produces a playable social simulation world with places, characters, items, rules, visual assets, and live human interaction.

The technical thesis is stricter:

\begin{itemize}
  \item Natural language is only the input surface.
  \item JSON remains the authoring contract.
  \item SQLite \texttt{world\_package.db} remains the runtime bundle.
  \item Pixel UI and live mode consume package-backed state, not a hidden frontend model.
  \item Generated worlds become public only after package, asset, browser, and visual QA gates pass.
\end{itemize}

\section{System Scope}

Agora 2.0 owns the runnable integration layer:

\begin{longtable}{>{\raggedright\arraybackslash}p{0.26\linewidth}>{\raggedright\arraybackslash}p{0.66\linewidth}}
\textbf{Area} & \textbf{Scope} \\
\hline
Creator workflow & Browser and API workflow for draft creation, revision, art generation, QA, and publish. \\
World compiler & LLM-generated builder spec, registry resolution, compiled \texttt{world\_config.json}, scenario materialization, and package creation. \\
Asset pipeline & Generated map, agent atlases, manifest feeds, live-ready bootstrap data, and strict \texttt{PIXEL READ}. \\
Package runtime & \texttt{world\_package.db}, \texttt{package\_meta.json}, materialized static workspace, access-code export, and pull by code. \\
Pixel catalog & Public latest-by-seed world records with validation gating. \\
Live runtime & DB-backed live sessions, REST state/actions, WebSocket movement input, authoritative deltas, and write-behind spatial persistence. \\
Operational validation & Backend startup smoke test, headless Firefox Pixel launch validation, Gemini MAP QA, regression scripts, and load tests. \\
\end{longtable}

\section{End-to-End Flow}

\subsection{Draft Creation}

The creator UI posts to \texttt{/api/world-builder/drafts} with a world name, genre, player count target, agent count target, focus, seed, and brief. The backend creates a durable draft directory under \texttt{output/world\_creator\_drafts/}, writes a manifest, creates revision \texttt{r001}, writes a placeholder status, queues a systemd user worker, and returns immediately.

This is a core design constraint. Long generation must not block the HTTP request path.

\subsection{Generation Worker}

The queued command is:

\begin{verbatim}
python -m agora_ui.world_builder generate-worker \
  --package-root ... --draft-id ... --revision-id ...
\end{verbatim}

The generation worker uses \texttt{VertexJsonClient} and the world-builder node pipeline to create planner, rooms, items, roles, main characters, hooks, visual direction, and gameplay loops. The normalized builder spec is compiled through \texttt{build\_world\_pipeline()}, which resolves registry-backed kits and policies into concrete artifacts.

The worker emits the revision artifact set:

\begin{itemize}
  \item \texttt{builder\_spec.json}
  \item \texttt{planner.json}
  \item \texttt{rooms\_spec.json}
  \item \texttt{items\_spec.json}
  \item \texttt{agents\_spec.json}
  \item \texttt{pixel\_frontend\_spec.json}
  \item \texttt{compiler\_report.json}
  \item \texttt{compiler\_critique.json}
  \item \texttt{world\_config.json}
  \item \texttt{scenario/}
  \item \texttt{world\_summary.md}
  \item \texttt{world\_package.db}
  \item \texttt{status.json}
\end{itemize}

\subsection{Art and QA Worker}

The current art worker executes three major commands:

\begin{enumerate}
  \item \texttt{python -m macro\_ui.build\_macro\_ui}
  \item \texttt{asset\_pipeline/generate\_world\_asset\_set.py}
  \item \texttt{asset\_pipeline/build\_live\_ready\_feed.py}
\end{enumerate}

It then repacks the revision package with only the assets matching the current \texttt{world\_id} and \texttt{world\_revision}. The generated map is injected into:

\begin{itemize}
  \item \texttt{world\_config} under the Pixel frontend map asset field
  \item \texttt{scenario/map\_grid.json} under \texttt{map\_visual.background\_url}
\end{itemize}

This prevents the package from passing with assets that exist on disk but are not connected to the runtime render path.

\subsection{Validation Gates}

A revision becomes publish-ready only after:

\begin{itemize}
  \item \texttt{PIXEL READ} passes for the isolated package workspace.
  \item Backend startup validation passes.
  \item Pixel UI launches in headless Firefox.
  \item The temporary validation package clears the \texttt{validation\_probe} lifecycle.
  \item Gemini MAP QA accepts both the stitched map asset and browser screenshot.
\end{itemize}

If \texttt{PIXEL READ} fails once, the live-ready feed may be retried once. Other critical failures must hardfail rather than silently producing placeholders.

\subsection{Publish}

Publishing exports the revision-local \texttt{world\_package.db}; it must not rebuild from raw config. The public package lands under:

\begin{verbatim}
output/package_exports/{access_code}/
  world_package.db
  package_meta.json
  materialized/
\end{verbatim}

The exported package receives a 16-character access code and is visible to Pixel only if it passes public catalog rules.

\section{Feature Inventory}

\begin{longtable}{>{\raggedright\arraybackslash}p{0.3\linewidth}>{\raggedright\arraybackslash}p{0.6\linewidth}}
\textbf{Feature} & \textbf{Current behavior} \\
\hline
One-brief generation & Converts a natural-language brief into a normalized builder spec and compiled runtime config. \\
Draft/revision history & Stores every revision with status, summary, package, and compiler artifacts. \\
Feedback revision & Creates a new revision from user feedback and prior summary context. \\
Registry-backed world assembly & Uses profiles, component kits, frontend affordances, item collections, inventory layers, property policies, and knowledge policies. \\
Protagonist-first worlds & Supports main-character-led worlds without requiring anonymous regular role groups. \\
Dynamic item catalog & Allows generated item catalogs while self-healing mandatory frontend affordance items. \\
Room-scoped visual metadata & Keeps room purpose, biome, decor, floor tile, wall tile, palette, and dimensions available to map and asset stages. \\
Package DB runtime & Stores files, metadata, and structured definitions inside \texttt{world\_package.db}. \\
Public catalog gating & Filters by \texttt{PIXEL READ}, startup status, source label, probe flag, and latest-by-seed deduplication. \\
Live multiplayer & Uses REST for boot/actions/state and WebSocket for realtime movement with authoritative deltas. \\
\end{longtable}

\section{Boundary Conditions}

Agora 2.0 is intentionally not:

\begin{itemize}
  \item A frontend-only sandbox.
  \item A JSON-only live runtime.
  \item A package-optional demo.
  \item A system that accepts placeholder art as success.
  \item A runtime where the browser owns inventory, trade, or authoritative position state.
\end{itemize}

The central boundary is that generated content may be creative, but runtime truth must stay explicit, materialized, validated, and package-backed.

\section{Operational Requirements}

\begin{itemize}
  \item The main runtime uses the \texttt{new\_py310} conda environment or the Python chosen by \texttt{resolve\_runtime\_python()}.
  \item Creator workers load \texttt{\$HOME/.config/agora\_ui\_runtime.env}.
  \item Fresh creator regression expects \texttt{AGORA\_AISTUDIO\_API\_KEY} and \texttt{AGORA\_VERTEX\_API\_KEY}.
  \item FLUX must be launched through \texttt{scripts/launch\_flux\_asset\_service.sh} so CUDA library paths are configured.
  \item Pixel assets should be served from the materialized package path for fast map and atlas bootstrap.
\end{itemize}

\section{Risk Register}

\begin{longtable}{>{\raggedright\arraybackslash}p{0.32\linewidth}>{\raggedright\arraybackslash}p{0.58\linewidth}}
\textbf{Risk} & \textbf{Mitigation} \\
\hline
Long generation blocks the server & Keep generation and art in detached systemd user workers with pollable status files. \\
Stale global assets leak into a new package & Repack from an isolated revision workspace and filter assets by world ID and revision. \\
World appears before validation completes & Preserve the \texttt{validation\_probe} gate and clear it only after successful headless launch. \\
Duplicate generated revisions crowd the catalog & Deduplicate server-side by seed, then world ID, then world name. \\
Map exists but is not rendered & Inject generated map paths into both config and scenario map grid before repack. \\
Pixel boot stalls on asset hydration & Prefer materialized static package assets over slow dynamic file serving. \\
FLUX silently fails and placeholders slip through & Start FLUX with the CUDA-aware wrapper and hardfail critical asset errors. \\
\end{longtable}

\section{Conclusion}

Agora 2.0's feature scope is best understood as a complete world compiler and runtime, not merely a prompt-to-image or prompt-to-JSON tool. The system earns the slogan ``one sentence, one world'' by preserving a chain of evidence from brief, to builder spec, to compiled JSON, to package DB, to asset manifests, to browser validation, to live DB-backed play.

\end{document}
